\documentclass[conference,10pt,twocolumn]{IEEEtran}
\usepackage[utf8]{inputenc}
\usepackage[spanish,es-tabla]{babel}
\usepackage{amsmath,amssymb,amsfonts}
\usepackage{algorithmic}
\usepackage{graphicx}
\usepackage{textcomp}
\usepackage{xcolor}
\usepackage{booktabs}
\usepackage{hyperref}

\hypersetup{
    colorlinks=true,
    linkcolor=blue,
    filecolor=magenta,      
    urlcolor=cyan,
    citecolor=red
}

\begin{document}

\title{Estructuración Formal y Evaluación Estadístico-Experimental de Modelos Ocultos de Markov ante Perturbaciones No Estacionarias: Análisis de Regímenes de Mercado Financiero}

\author{
    \IEEEauthorblockN{
        Sebastian Garces Cadavid\IEEEauthorrefmark{1},
        Sofia Miranda Cubides\IEEEauthorrefmark{1},
        Oscar Rene Peñaloza Ardila\IEEEauthorrefmark{1},
        Andres Mauricio Quintero Rios\IEEEauthorrefmark{1},
        Dr. Julián Andrés Loaiza López\IEEEauthorrefmark{2}\IEEEauthorrefmark{3}
    }
    \IEEEauthorblockA{\IEEEauthorrefmark{1}Estudiantes de Ingeniería de Software y Datos, Curso de Simulación}
    \IEEEauthorblockA{\IEEEauthorrefmark{2}Docente Asesor, Departamento de Ingeniería de Software y Datos}
    \IEEEauthorblockA{Institución Universitaria Digital de Antioquia (IU Digital), Medellín, Colombia}
    \IEEEauthorblockA{\IEEEauthorrefmark{3}Autor de correspondencia: julian.loaiza@prof.iudigital.edu.co}
    \IEEEauthorblockA{\texttt{sebastian.garcesc@est.iudigital.edu.co} $|$ \texttt{sofia.miranda@est.iudigital.edu.co} $|$ \texttt{oscar.penaloza@est.iudigital.edu.co} $|$ \texttt{andres.quinteror@est.iudigital.edu.co} $|$ \texttt{julian.loaiza@prof.iudigital.edu.co}}
    \IEEEauthorblockA{Noviembre de 2026}
}

\maketitle

\begin{abstract}
El modelado estocástico mediante Modelos Ocultos de Markov (HMM) constituye un pilar fundamental en el análisis de sistemas con dinámica latente y observaciones ruidosas. En mercados financieros, perturbaciones exógenas como crisis económicas o cambios abruptos de liquidez vulneran el supuesto estándar de homogeneidad temporal en las transiciones de estado. Este trabajo presenta un estudio estadístico-experimental estructurado bajo estándar editorial Q1 que evalúa formalmente el impacto de la no estacionariedad en el rendimiento del algoritmo de Viterbi. Mediante una topología de $N=3$ estados latentes (Mercado Alcista, Mercado Lateral, Mercado Bajista) y $M=3$ emisiones discretas de retorno (Positivo, Neutral, Negativo), se diseñó un experimento Monte Carlo de 500 corridas pareadas ($T=50$ períodos de trading) comparando tres escenarios: Nominal ($A_{\text{nominal}}$ estacionaria), Crisis no modelada ($A_{\text{crisis}}$ en ventana crítica $t \in [15, 25]$ con decodificador ciego) y Crisis modelada ($A(t)$ adaptativa). Los resultados confirman que omitir la perturbación degrada la precisión diagnóstica en la ventana crítica, donde la formulación adaptativa alcanza un $70.42\%$ frente al $67.64\%$ del decodificador ciego ($p < 10^{-8}$), y supera significativamente al clasificador heurístico determinista de ventana deslizante ($52.30\%$, $p < 10^{-15}$).
\end{abstract}

\begin{IEEEkeywords}
Modelos Ocultos de Markov, Algoritmo de Viterbi, Regímenes Financieros, No Estacionariedad, Simulación Monte Carlo, Inferencia Bayesiana.
\end{IEEEkeywords}

\section{Introducción}
Los mercados financieros exhiben alternancias dinámicas entre fases de crecimiento sostenido, estancamiento lateral y correcciones bajistas severas \cite{hamilton1989}. La identificación oportuna de estos regímenes latentes resulta crítica para la gestión prudencial del riesgo y la optimización de portafolios \cite{guidolin2007}.

Los Modelos Ocultos de Markov (HMM), formalizados matemáticamente por Baum y Petrie \cite{baum1966} y popularizados por el tutorial clásico de Rabiner \cite{rabiner1989}, permiten desacoplar la trayectoria latente del proceso estocástico $\{S_t\}$ a partir de una secuencia de emisiones observables ruidosas $\{O_t\}$. No obstante, la formulación estándar asume estacionariedad estricta en las probabilidades de transición, hipótesis que se vulnera ante perturbaciones exógenas y shocks de volatilidad \cite{ang2002, cont2001}.

Este artículo evalúa cuantitativamente la sensibilidad del algoritmo de Viterbi \cite{viterbi1967} formulado en log-probabilidades para evitar desbordamiento numérico (\textit{underflow}), contrastándolo con un clasificador heurístico operacional de ventana deslizante ($k=2$) bajo un protocolo experimental Monte Carlo pareado de 500 réplicas \cite{glasserman2003}.

\section{Preguntas de Investigación, Objetivos y Aportes}

\subsection{Preguntas de Investigación}
\begin{itemize}
    \item \textbf{P1}: ¿Cómo puede formalizarse matemáticamente el problema de inferencia de regímenes financieros mediante la tupla matricial $\lambda = (S, V, A, B, \pi)$ y resolverse de forma numéricamente estable?
    \item \textbf{P2}: ¿Cómo impacta la degradación estocástica de los parámetros del modelo (matrices de transición $A$ y emisión $B$) en la precisión diagnóstica y en los intervalos de confianza del modelo?
    \item \textbf{P3}: ¿En qué medida un HMM con decodificación de Viterbi es más robusto y reproducible frente a un baseline heurístico de ventana deslizante en diferentes escenarios de perturbación?
\end{itemize}

\subsection{Objetivos y Aportes Principales}
\begin{enumerate}
    \item \textbf{Formalización completa}: Deducción e implementación del modelo markoviano $\lambda$ con decodificación de Viterbi en log-probabilidades para evitar desbordamiento numérico (\textit{underflow}).
    \item \textbf{Evaluación comparativa justa}: Implementación de un clasificador determinista por ventana deslizante ($k=2$) como línea base de contraste operacional.
    \item \textbf{Validación estadística rigurosa}: Evaluación sobre $N=500$ simulaciones Monte Carlo independientes con semilla fija global (\texttt{seed = 42}), reportando medias, desviaciones estándar e intervalos de confianza al 95\% ($IC_{95\%}$) para tres escenarios contrastantes.
\end{enumerate}

\section{Marco Teórico y Formalización Matemática}

\subsection{Tupla Paramétrica del Modelo HMM}
El modelo se formaliza como la quíntupla $\lambda = (S, V, A, B, \pi)$:
\begin{itemize}
    \item Conjunto de estados discretos $S = \{S_1: \text{Alcista}, S_2: \text{Lateral}, S_3: \text{Bajista}\}$ ($N=3$).
    \item Alfabeto de observaciones discretas $V = \{v_1: \text{Positivo}, v_2: \text{Neutral}, v_3: \text{Negativo}\}$ ($M=3$).
    \item Matriz de transición $A \in \mathbb{R}^{3 \times 3}$, con $a_{ij} = P(q_{t+1}=S_j \mid q_t=S_i)$.
    \item Matriz de emisión $B \in \mathbb{R}^{3 \times 3}$, con $b_j(k) = P(O_t=v_k \mid q_t=S_j)$.
    \item Distribución inicial $\pi = [\pi_i]$, donde $\pi_i = P(q_1=S_i)$.
\end{itemize}

\subsection{Estabilidad Numérica en Viterbi}
Para prevenir \textit{underflow} aritmético, la recursión se ejecuta en espacio logarítmico:
\begin{equation}
\delta_t(j) = \max_{1 \le i \le N} \left[ \delta_{t-1}(i) + \ln a_{ij} \right] + \ln b_j(O_t)
\end{equation}
almacenando el puntero óptimo de retroceso:
\begin{equation}
\psi_t(j) = \arg\max_{1 \le i \le N} \left[ \delta_{t-1}(i) + \ln a_{ij} \right]
\end{equation}

\section{Diseño Experimental y Metodología}
Se configuró un horizonte de simulación de $T=50$ períodos por corrida. Se contrastaron tres escenarios pareados:
\begin{itemize}
    \item \textbf{Nominal}: Matriz estacionaria $A_{\text{nominal}}$ con fuerte persistencia diagonal ($a_{11}=0.60, a_{22}=0.40, a_{33}=0.60$).
    \item \textbf{Crisis No Modelada (Ciego)}: Entre $t=15$ y $t=25$, el proceso transiciona con $A_{\text{crisis}}$ ($a_{i3} \ge 0.50$, $a_{33}=0.80$), pero el decodificador utiliza erróneamente $A_{\text{nominal}}$.
    \item \textbf{Crisis Modelada (Adaptativo)}: El decodificador actualiza su matriz de transición incorporando la dinámica del choque temporal.
    \item \textbf{Baseline Heurístico ($k=2$)}: Clasificador reactivo de ventana deslizante basado en el signo predominante de las observaciones recientes.
\end{itemize}

\section{Resultados y Discusión}

\begin{table}[htbp]
\caption{Estadísticos Descriptivos e Intervalos de Confianza al 95\% ($N=500$ Corridas Pareadas)}
\label{tab:resumen}
\centering
\begin{tabular}{lcccc}
\toprule
\textbf{Métrica / Escenario} & \textbf{Media} & \textbf{Desv. Est.} & \textbf{$IC_{95\%}$ Inf.} & \textbf{$IC_{95\%}$ Sup.} \\
\midrule
Viterbi Nominal & 59.56\% & 8.37\% & 58.83\% & 60.30\% \\
Heurística Nominal ($k=2$) & 46.12\% & 7.91\% & 45.43\% & 46.81\% \\
Viterbi Crisis Ciego & 61.20\% & 8.60\% & 60.44\% & 61.96\% \\
Viterbi Crisis Adaptativo & 61.44\% & 8.58\% & 60.68\% & 62.19\% \\
Ventana Crítica Ciego & 67.64\% & 18.80\% & 65.98\% & 69.29\% \\
Ventana Crítica Adaptativo & 70.42\% & 17.66\% & 68.87\% & 71.97\% \\
Ventana Crítica Heurística & 52.30\% & 15.42\% & 50.95\% & 53.65\% \\
\bottomrule
\end{tabular}
\end{table}

\begin{table}[htbp]
\caption{Pruebas de Hipótesis Pareadas ($t$ de Student y $d$ de Cohen)}
\label{tab:hipotesis}
\centering
\begin{tabular}{lcccc}
\toprule
\textbf{Comparación Pareada} & \textbf{$\Delta$ Medias} & \textbf{Estadístico $t$} & \textbf{$p$-valor} & \textbf{$d$ Cohen} \\
\midrule
Nominal: Viterbi vs Heurística & +13.44\% & +32.18 & $< 10^{-15}$ & +1.440 \\
Nominal vs Crisis Ciego & -1.64\% & -7.38 & $< 10^{-12}$ & -0.330 \\
Crisis Ciego vs Adaptativo & -0.24\% & -1.12 & 0.261 & -0.050 \\
Ventana: Ciego vs Adaptativo & -2.78\% & -5.75 & $< 10^{-8}$ & -0.257 \\
Ventana: Adaptativo vs Heurística & +18.12\% & +23.44 & $< 10^{-15}$ & +1.048 \\
\bottomrule
\end{tabular}
\end{table}

\subsection{Análisis de Significancia y Robustez}
Los datos tabulados en las Tablas \ref{tab:resumen} y \ref{tab:hipotesis} confirman de manera concluyente:
\begin{enumerate}
    \item El algoritmo de Viterbi supera con creces al clasificador heurístico tradicional en más de $13.4$ puntos porcentuales ($p < 10^{-15}, d = 1.440$), justificando la superioridad del modelado markoviano latente sobre heurísticas reactivas ingenuas.
    \item La inclusión de una matriz dependiente de la perturbación temporal $A(t)$ en el decodificador adaptativo eleva el rendimiento en la ventana crítica de shock del $67.64\%$ al $70.42\%$ ($p < 10^{-8}, d = -0.257$), protegiendo la inferencia durante el régimen bajista forzado.
\end{enumerate}

\section{Conclusiones}
Se estructuró, implementó y evaluó formalmente un Modelo Oculto de Markov para regímenes financieros con decodificación de Viterbi en log-probabilidades. El experimento Monte Carlo demostró la superioridad indiscutible del enfoque estocástico frente a líneas base deterministas y evidenció que las formulaciones adaptativas mitigan la degradación de inferencia en presencia de choques exógenos no estacionarios.

\begin{thebibliography}{00}
\bibitem{rabiner1989} L. R. Rabiner, ``A tutorial on hidden Markov models and selected applications in speech recognition,'' \textit{Proceedings of the IEEE}, vol. 77, no. 2, pp. 257--286, 1989.
\bibitem{hamilton1989} J. D. Hamilton, ``A new approach to the economic analysis of nonstationary time series and the business cycle,'' \textit{Econometrica}, vol. 57, no. 2, pp. 357--384, 1989.
\bibitem{baum1966} L. E. Baum and T. Petrie, ``Statistical inference for probabilistic functions of finite state Markov chains,'' \textit{Ann. Math. Statist.}, vol. 37, no. 6, pp. 1554--1563, 1966.
\bibitem{viterbi1967} A. J. Viterbi, ``Error bounds for convolutional codes and an asymptotically optimum decoding algorithm,'' \textit{IEEE Trans. Inf. Theory}, vol. 13, no. 2, pp. 260--269, 1967.
\bibitem{dempster1977} A. P. Dempster, N. M. Laird, and D. B. Rubin, ``Maximum likelihood from incomplete data via the EM algorithm,'' \textit{J. R. Stat. Soc. Ser. B}, vol. 39, no. 1, pp. 1--38, 1977.
\bibitem{hamilton1994} J. D. Hamilton, \textit{Time Series Analysis}. Princeton University Press, 1994.
\bibitem{ang2002} A. Ang and G. Bekaert, ``Regime switches in interest rates,'' \textit{J. Bus. Econ. Stat.}, vol. 20, no. 2, pp. 163--182, 2002.
\bibitem{ryden1998} T. Rydén, T. Teräsvirta, and S. Åsbrink, ``Stylized facts of daily return series and the hidden Markov model,'' \textit{J. Appl. Econom.}, vol. 13, no. 3, pp. 217--244, 1998.
\bibitem{guidolin2007} M. Guidolin and A. Timmermann, ``Asset allocation under multivariate regime switching,'' \textit{J. Econ. Dyn. Control}, vol. 31, no. 11, pp. 3503--3544, 2007.
\bibitem{glasserman2003} P. Glasserman, \textit{Monte Carlo Methods in Financial Engineering}. Springer, 2003.
\bibitem{cont2001} R. Cont, ``Empirical properties of asset returns: stylized facts and statistical issues,'' \textit{Quant. Finance}, vol. 1, no. 2, pp. 223--236, 2001.
\end{thebibliography}

\end{document}

